% Topic presentation -- 20 minutes.
% No animations and no screenshots of papers: equations are typeset in LaTeX.
\documentclass[aspectratio=169,11pt]{beamer}
\usepackage[utf8]{inputenc}
\usepackage[T1]{fontenc}
\usepackage[english]{babel}
\usepackage{amsmath,amssymb}
\usepackage{booktabs}
\usetheme{default}
\usecolortheme{seahorse}
\setbeamertemplate{navigation symbols}{}
\setbeamertemplate{footline}[frame number]
\setbeamerfont{frametitle}{size=\large}
\setbeamercolor{alerted text}{fg=blue!55!black}

\newcommand{\Rb}{R_B}
\newcommand{\Rm}{R_M}
\newcommand{\kb}{k_B}
\newcommand{\km}{k_M}

\title{Credit Quantity, Credit Prices, and Future Output}
\subtitle{A model of heterogeneous financing access -- Track B}
\author{Alvaro Marcelo Ch\'avez Unyen}
\date{October 7, 2026}

\begin{document}

\begin{frame}[plain]
  \titlepage
  \begin{center}\small
    \url{https://github.com/amchavezu/ai-project}
  \end{center}
\end{frame}

\begin{frame}{1 \textperiodcentered{} The question and Track B}
\large
When does the \alert{quantity} of bank credit contain more information about
future production than the \alert{price} of credit?

\medskip
\normalsize
\begin{itemize}
  \item Track B: formulate the firm model missing from my undergraduate thesis.
  \item Preliminary estimates suggest different credit indicators may be useful
        for non-tradable and tradable activity.
  \item I treat that evidence as qualitative and unfinished. The model targets
        the cross-sector mechanism, not exact horizons or causality.
\end{itemize}
\end{frame}

\begin{frame}{1 \textperiodcentered{} The pattern that needs a mechanism}
\begin{columns}[T,onlytextwidth]
  \column{.47\linewidth}
  \textbf{Preliminary quantity pattern}

  Bank credit volume appears more informative for non-tradable activity at a
  relatively short horizon.

  \medskip
  Candidate mechanism: bank-dependent firms cannot replace a missing loan.

  \column{.47\linewidth}
  \textbf{Preliminary price pattern}

  A credit-price indicator may become more informative for tradable activity
  over a longer horizon.

  \medskip
  Candidate mechanism: firms with alternative finance choose scale using the
  marginal funding cost.
\end{columns}

\vfill
\footnotesize These are working hypotheses from a thesis in progress, not
final empirical findings.
\end{frame}

\begin{frame}{2 \textperiodcentered{} The model missing from the thesis}
At date zero, firm $i$ has internal funds $n_i$ and chooses capital $k_i$,
bank credit $b_i$, and outside finance $m_i$:
\[
 \max_{k_i,b_i,m_i}\quad
 \beta F_i(k_i)-\Rb b_i-\Rm m_i
\]
\[
 k_i=n_i+b_i+m_i,\qquad
 0\leq b_i\leq\bar b_i,\qquad m_i\geq0.
\]

\begin{itemize}
  \item $F_i'(k)>0$ and $F_i''(k)<0$.
  \item $0<\Rb\leq\Rm=\Rb+\tau_i$: bank funds are cheaper but capped.
  \item A bank-dependent firm has $m_i=0$; a market-access firm may have
        $m_i>0$.
\end{itemize}
\end{frame}

\begin{frame}{2 \textperiodcentered{} Two margins of transmission}
Let $\bar k_i=n_i+\bar b_i$ be capital after exhausting the bank line.

\medskip
\begin{center}
\begin{tabular}{p{.24\linewidth}p{.31\linewidth}p{.31\linewidth}}
\toprule
 & \textbf{Quantity-sensitive kink} & \textbf{Price-sensitive margin}\\
\midrule
Marginal finance & No cheap substitute for the bank line & Outside finance\\
Choice of capital & $k_i^*=\bar k_i$ & $\beta F_i'(k_i^*)=\Rm$\\
Larger bank line & Raises total capital & Replaces outside finance\\
Higher common spread & No local scale effect at the kink & Lowers desired capital\\
\bottomrule
\end{tabular}
\end{center}

\medskip
The comparison is local: small changes must leave the firm in the same regime.
\end{frame}

\begin{frame}{3 \textperiodcentered{} The result I expect}
With $F_i(k)=a_i k-c_i k^2/2$, define
\[
 \kb=\frac{\beta a_i-\Rb}{\beta c_i},\qquad
 \km=\frac{\beta a_i-\Rm}{\beta c_i},\qquad \km\leq\kb.
\]
\[
k_i^*=\begin{cases}
n_i, & \kb\leq n_i,\\
\kb, & n_i<\kb<\bar k_i,\\
\bar k_i, & \km\leq\bar k_i\leq\kb,\\
\km, & \bar k_i<\km.
\end{cases}
\]

\small
The third line is the binding quantity regime. The fourth is the
market-finance regime. Strict inequalities deliver the comparative statics
without a regime switch.
\end{frame}

\begin{frame}{3 \textperiodcentered{} Expected sectoral ranking}
Let $\lambda_s$ be the share of sector $s$ at the quantity-sensitive kink:
\[
Y_s=\lambda_s F(n+\bar b)+(1-\lambda_s)F(k_M).
\]
If $\lambda_N>\lambda_T$ and marginal products are positive, I expect
\[
 \frac{\partial Y_N}{\partial\bar b}>
 \frac{\partial Y_T}{\partial\bar b},
 \qquad
 \left|\frac{\partial Y_N}{\partial z}\right|<
 \left|\frac{\partial Y_T}{\partial z}\right|.
\]

\begin{itemize}
  \item Credit quantity signals the feasible project scale for constrained firms.
  \item The spread signals the marginal cost and desired project scale for firms
        that can substitute across funding sources.
  \item Adjustment speeds may affect timing, but the baseline does not derive
        exact forecast horizons.
\end{itemize}
\end{frame}

\begin{frame}{4 \textperiodcentered{} Why this is not already solved}
\small
\begin{tabular}{p{.25\linewidth}p{.30\linewidth}p{.35\linewidth}}
\toprule
\textbf{Closest work} & \textbf{What it establishes} & \textbf{What remains for this project}\\
\midrule
Holmstr\"om--Tirole (1997) & Capital, monitoring, and access to direct finance & No sectoral ranking of quantity and price predictors\\
Gilchrist--Zakraj\v{s}ek (2012) & Credit spreads predict aggregate activity & No binding bank-quantity comparison at the firm level\\
M\"uller--Verner (2024) & Sectoral credit allocation predicts different macro outcomes & No capped-bank versus marginal-market threshold\\
\bottomrule
\end{tabular}

\medskip
I checked the published versions, public working-paper versions, and the
relevant appendices and follow-on investment paper. The contribution is a
narrow synthesis tailored to the thesis, not a claim that the underlying
financial-friction mechanisms are new.
\end{frame}

\begin{frame}{5 \textperiodcentered{} Plan and risks}
\begin{columns}[T,onlytextwidth]
  \column{.48\linewidth}
  \textbf{Plan}
  \begin{itemize}
    \item Prove the piecewise optimum and sectoral rankings.
    \item Check every reported result with \texttt{code/verify.py}.
    \item Simulate behavior near regime boundaries.
    \item Formalize every numbered result through AppliedModelingLib.
  \end{itemize}

  \column{.48\linewidth}
  \textbf{Main risks}
  \begin{itemize}
    \item Outside-finance access is not directly measured in the thesis data.
    \item Endogenous credit limits may remove the sharp local zero.
    \item A static model cannot identify exact forecast timing.
    \item The piecewise global optimum is the hardest Lean step.
  \end{itemize}
\end{columns}

\vfill
\small If a step fails, I will keep the result conditional and report the exact
open case rather than strengthen the assumptions silently.
\end{frame}

\end{document}
